
\documentclass[11pt, a4paper]{article}
\usepackage[margin=1in]{geometry}
\usepackage{amsmath, amssymb}
\usepackage{xcolor}
\usepackage{booktabs}
\usepackage{parskip}
\usepackage{fancyhdr}
\usepackage{tcolorbox}
\setlength{\headheight}{14pt}

\definecolor{secblue}{RGB}{26, 71, 111}
\definecolor{notegreen}{RGB}{40, 120, 60}
\definecolor{explgray}{RGB}{100, 100, 100}
\definecolor{lightblue}{RGB}{230, 242, 255}
\definecolor{lightyellow}{RGB}{255, 250, 230}

\newcommand{\gut}[1]{\begin{tcolorbox}[colback=lightyellow, colframe=orange!50, title=Gut feeling]#1\end{tcolorbox}}
\newcommand{\recap}[1]{\begin{tcolorbox}[colback=lightblue, colframe=secblue!50, title=Recap]#1\end{tcolorbox}}

\pagestyle{fancy}
\fancyhf{}
\rhead{\textit{OLS Proof \& Worked Example}}
\lhead{\textit{Arif's GP Study}}
\cfoot{\thepage}

\title{\color{secblue}\textbf{Ordinary Least Squares}\\[0.3em]
\Large Proof of $\mathbf{w} = (X^\top X)^{-1} X^\top \mathbf{y}$ and Full Numerical Example}
\author{}
\date{}

\begin{document}
\maketitle
\thispagestyle{fancy}


% ============================================================
\section*{\color{secblue}Part 1: The Setup}

We have $n$ data points. Each data point has an input $x_i$ and an output $y_i$.
We want to fit a line:

\begin{equation}
\hat{y}_i = w_0 + w_1 \, x_i
\tag{OLS-1}
\end{equation}

\begin{center}
\begin{tabular}{cl}
\toprule
\textbf{Symbol} & \textbf{Plain English} \\
\midrule
$\hat{y}_i$ & Our prediction for sample $i$ \\
$y_i$ & The actual observed value for sample $i$ \\
$w_0$ & The intercept (value when $x = 0$) \\
$w_1$ & The slope (how much $y$ changes per unit of $x$) \\
$x_i$ & The input feature of sample $i$ \\
\bottomrule
\end{tabular}
\end{center}

In matrix form, we stack all $n$ predictions into one equation:

\begin{equation}
\hat{\mathbf{y}} = X \, \mathbf{w}
\tag{OLS-2}
\end{equation}

where:

\begin{equation}
X = \begin{pmatrix} 1 & x_1 \\ 1 & x_2 \\ \vdots & \vdots \\ 1 & x_n \end{pmatrix}
\quad
\mathbf{w} = \begin{pmatrix} w_0 \\ w_1 \end{pmatrix}
\quad
\hat{\mathbf{y}} = \begin{pmatrix} w_0 + w_1 x_1 \\ w_0 + w_1 x_2 \\ \vdots \\ w_0 + w_1 x_n \end{pmatrix}
\tag{OLS-3}
\end{equation}

\gut{
The column of ones in $X$ is a trick. When you multiply a row $[1, \; x_i]$ by $[w_0, \; w_1]^\top$, you get $1 \cdot w_0 + x_i \cdot w_1 = w_0 + w_1 x_i$. That's how the intercept sneaks into the matrix form.
}


% ============================================================
\section*{\color{secblue}Part 2: The Goal --- Minimize Squared Errors}

The \textbf{residual} for sample $i$ is how wrong our prediction is:

\begin{equation}
r_i = y_i - \hat{y}_i = y_i - (w_0 + w_1 x_i)
\tag{OLS-4}
\end{equation}

We want to minimize the \textbf{sum of squared residuals}:

\begin{equation}
L(\mathbf{w}) = \sum_{i=1}^{n} r_i^2 = \sum_{i=1}^{n} (y_i - \hat{y}_i)^2
\tag{OLS-5}
\end{equation}

\begin{center}
$L$ = ``loss function'' --- the thing we want to make as small as possible.
\end{center}

In matrix form, the vector of all residuals is:

\begin{equation}
\mathbf{r} = \mathbf{y} - X\mathbf{w}
\tag{OLS-6}
\end{equation}

And the sum of squared residuals is:

\begin{equation}
L(\mathbf{w}) = \mathbf{r}^\top \mathbf{r} = (\mathbf{y} - X\mathbf{w})^\top (\mathbf{y} - X\mathbf{w})
\tag{OLS-7}
\end{equation}

\gut{
$\mathbf{r}^\top \mathbf{r}$ is just the dot product of the residual vector with itself, which equals $r_1^2 + r_2^2 + \cdots + r_n^2$. It's a single number that measures total error.
}


% ============================================================
\section*{\color{secblue}Part 3: Expanding the Loss Function}

Let's expand Eq.~OLS-7 step by step. Using the rule $(A - B)^\top = A^\top - B^\top$:

\begin{equation}
L = (\mathbf{y}^\top - \mathbf{w}^\top X^\top)(\mathbf{y} - X\mathbf{w})
\tag{OLS-8}
\end{equation}

Multiply out (just like $(a - b)(c - d) = ac - ad - bc + bd$):

\begin{align}
L &= \mathbf{y}^\top \mathbf{y} \;-\; \mathbf{y}^\top X\mathbf{w} \;-\; \mathbf{w}^\top X^\top \mathbf{y} \;+\; \mathbf{w}^\top X^\top X \mathbf{w}
\tag{OLS-9}
\end{align}

\begin{center}
\begin{tabular}{cl}
\toprule
\textbf{Term} & \textbf{Plain English} \\
\midrule
$\mathbf{y}^\top \mathbf{y}$ & Total energy in the data (a constant, doesn't depend on $\mathbf{w}$) \\
$\mathbf{y}^\top X\mathbf{w}$ & How well our prediction aligns with the data \\
$\mathbf{w}^\top X^\top \mathbf{y}$ & Same as above (it's a scalar, so it equals its own transpose) \\
$\mathbf{w}^\top X^\top X \mathbf{w}$ & Total energy in our predictions \\
\bottomrule
\end{tabular}
\end{center}

Since $\mathbf{y}^\top X \mathbf{w}$ is a scalar (a $1\times 1$ matrix), it equals its transpose: $\mathbf{w}^\top X^\top \mathbf{y}$. So the two middle terms are identical, and we can simplify:

\begin{equation}
L = \mathbf{y}^\top \mathbf{y} \;-\; 2\,\mathbf{w}^\top X^\top \mathbf{y} \;+\; \mathbf{w}^\top X^\top X \,\mathbf{w}
\tag{OLS-10}
\end{equation}


% ============================================================
\section*{\color{secblue}Part 4: Taking the Derivative and Setting It to Zero}

To find the minimum, we take the derivative of $L$ with respect to $\mathbf{w}$ and set it to zero.

\recap{
Two matrix calculus rules we need (analogous to scalar rules):
\begin{enumerate}
\item $\dfrac{\partial}{\partial \mathbf{w}} (\mathbf{a}^\top \mathbf{w}) = \mathbf{a}$
\hfill \textit{(like $\frac{d}{dx}(ax) = a$ in scalar calculus)}
\item $\dfrac{\partial}{\partial \mathbf{w}} (\mathbf{w}^\top A \,\mathbf{w}) = 2A\mathbf{w}$
\hfill \textit{(like $\frac{d}{dx}(ax^2) = 2ax$ in scalar calculus)}
\end{enumerate}
These are the matrix versions of the two most basic derivative rules you already know.
}

Apply these rules to each term of Eq.~OLS-10:

\begin{align}
\frac{\partial L}{\partial \mathbf{w}} &= 0 \;-\; 2\,X^\top \mathbf{y} \;+\; 2\,X^\top X\,\mathbf{w}
\tag{OLS-11}
\end{align}

\begin{center}
\begin{tabular}{cl}
\toprule
\textbf{Term in $L$} & \textbf{Its derivative w.r.t.\ $\mathbf{w}$} \\
\midrule
$\mathbf{y}^\top \mathbf{y}$ (constant) & $0$ \\
$-2\,\mathbf{w}^\top X^\top \mathbf{y}$ & $-2\,X^\top \mathbf{y}$ \quad (using rule 1 with $\mathbf{a} = X^\top \mathbf{y}$) \\
$\mathbf{w}^\top X^\top X\, \mathbf{w}$ & $2\,X^\top X\,\mathbf{w}$ \quad (using rule 2 with $A = X^\top X$) \\
\bottomrule
\end{tabular}
\end{center}

Set the derivative to zero:

\begin{equation}
-2\,X^\top \mathbf{y} + 2\,X^\top X\,\mathbf{w} = 0
\tag{OLS-12}
\end{equation}

Divide both sides by 2:

\begin{equation}
X^\top X\,\mathbf{w} = X^\top \mathbf{y}
\tag{OLS-13}
\end{equation}

\gut{
Eq.~OLS-13 is called the \textbf{normal equation}. It says: ``the optimal weights are the ones where $X^\top X \,\mathbf{w}$ exactly equals $X^\top \mathbf{y}$.'' This is a system of linear equations --- the matrix equivalent of $ax = b$.
}

Multiply both sides by $(X^\top X)^{-1}$ on the left:

\begin{equation}
\boxed{\mathbf{w} = (X^\top X)^{-1} \, X^\top \mathbf{y}}
\tag{OLS-14}
\end{equation}

That's it. This is the proof. $\square$

\gut{
Compare with scalar algebra: if $ax = b$, then $x = a^{-1}b = b/a$. Here, $X^\top X$ plays the role of $a$, $X^\top \mathbf{y}$ plays the role of $b$, and $\mathbf{w}$ plays the role of $x$. We ``divide'' both sides by $X^\top X$ (i.e., multiply by its inverse).
}


% ============================================================
% ============================================================
\newpage
\section*{\color{secblue}Part 5: Full Numerical Example with 2 Samples}

\subsection*{The Data}

\begin{center}
\begin{tabular}{ccc}
\toprule
\textbf{Sample} & \textbf{Size (sqm)} & \textbf{Price (lakh)} \\
\midrule
1 & 30 & 76.2 \\
2 & 80 & 201.7 \\
\bottomrule
\end{tabular}
\end{center}


\subsection*{Step 1: Build the Design Matrix $X$ and Target Vector $\mathbf{y}$}

\begin{equation}
X = \begin{pmatrix} 1 & 30 \\ 1 & 80 \end{pmatrix}
\quad
\mathbf{y} = \begin{pmatrix} 76.2 \\ 201.7 \end{pmatrix}
\tag{N-1}
\end{equation}

$X$ is $2 \times 2$. First column is all ones (for the intercept). Second column is the house sizes.


\subsection*{Step 2: Compute $X^\top$}

\begin{equation}
X^\top = \begin{pmatrix} 1 & 1 \\ 30 & 80 \end{pmatrix}
\tag{N-2}
\end{equation}

Just swap rows and columns of $X$.


\subsection*{Step 3: Compute $X^\top X$ \quad ($2\times 2 \;\cdot\; 2\times 2 = 2\times 2$)}

\begin{equation}
X^\top X = \begin{pmatrix} 1 & 1 \\ 30 & 80 \end{pmatrix}
\begin{pmatrix} 1 & 30 \\ 1 & 80 \end{pmatrix}
\tag{N-3}
\end{equation}

Element by element:

\begin{align}
(X^\top X)_{1,1} &= 1 \times 1 + 1 \times 1 = \mathbf{2} \tag{N-3a} \\
(X^\top X)_{1,2} &= 1 \times 30 + 1 \times 80 = \mathbf{110} \tag{N-3b} \\
(X^\top X)_{2,1} &= 30 \times 1 + 80 \times 1 = \mathbf{110} \tag{N-3c} \\
(X^\top X)_{2,2} &= 30 \times 30 + 80 \times 80 = 900 + 6400 = \mathbf{7300} \tag{N-3d}
\end{align}

\begin{equation}
X^\top X = \begin{pmatrix} 2 & 110 \\ 110 & 7300 \end{pmatrix}
\tag{N-4}
\end{equation}

\gut{
What do these numbers mean?
\begin{itemize}
\item $(X^\top X)_{1,1} = 2$ = number of samples (the ones column dotted with itself)
\item $(X^\top X)_{1,2} = 110$ = sum of all $x$ values ($30 + 80$)
\item $(X^\top X)_{2,2} = 7300$ = sum of squared $x$ values ($30^2 + 80^2$)
\end{itemize}
This matrix summarizes the ``geometry'' of your input data.
}


\subsection*{Step 4: Compute $X^\top \mathbf{y}$ \quad ($2\times 2 \;\cdot\; 2\times 1 = 2\times 1$)}

\begin{equation}
X^\top \mathbf{y} = \begin{pmatrix} 1 & 1 \\ 30 & 80 \end{pmatrix}
\begin{pmatrix} 76.2 \\ 201.7 \end{pmatrix}
\tag{N-5}
\end{equation}

\begin{align}
(X^\top \mathbf{y})_1 &= 1 \times 76.2 + 1 \times 201.7 = \mathbf{277.9} \tag{N-5a} \\
(X^\top \mathbf{y})_2 &= 30 \times 76.2 + 80 \times 201.7 = 2286 + 16136 = \mathbf{18422} \tag{N-5b}
\end{align}

\begin{equation}
X^\top \mathbf{y} = \begin{pmatrix} 277.9 \\ 18422 \end{pmatrix}
\tag{N-6}
\end{equation}

\gut{
\begin{itemize}
\item First entry = sum of all $y$ values ($76.2 + 201.7 = 277.9$)
\item Second entry = sum of $x_i \cdot y_i$ ($30 \times 76.2 + 80 \times 201.7 = 18422$)
\end{itemize}
This vector summarizes how inputs and outputs relate.
}


\subsection*{Step 5: Invert $X^\top X$}

Using our $2\times 2$ inverse formula (Eq.~11 from the matrix inverse derivation):

\begin{equation}
\det(X^\top X) = 2 \times 7300 - 110 \times 110 = 14600 - 12100 = 2500
\tag{N-7}
\end{equation}

\begin{equation}
(X^\top X)^{-1} = \frac{1}{2500} \begin{pmatrix} 7300 & -110 \\ -110 & 2 \end{pmatrix}
= \begin{pmatrix} 2.92 & -0.044 \\ -0.044 & 0.0008 \end{pmatrix}
\tag{N-8}
\end{equation}


\subsection*{Step 6: Compute $\mathbf{w} = (X^\top X)^{-1} \, X^\top \mathbf{y}$}

\begin{equation}
\mathbf{w} = \begin{pmatrix} 2.92 & -0.044 \\ -0.044 & 0.0008 \end{pmatrix}
\begin{pmatrix} 277.9 \\ 18422 \end{pmatrix}
\tag{N-9}
\end{equation}

\begin{align}
w_0 &= 2.92 \times 277.9 + (-0.044) \times 18422 = 811.47 - 810.57 = \mathbf{0.90} \tag{N-9a} \\
w_1 &= (-0.044) \times 277.9 + 0.0008 \times 18422 = -12.23 + 14.74 = \mathbf{2.51} \tag{N-9b}
\end{align}

\begin{equation}
\boxed{\mathbf{w} = \begin{pmatrix} 0.90 \\ 2.51 \end{pmatrix}}
\tag{N-10}
\end{equation}

So the fitted line is:

\begin{equation}
\hat{y} = 0.90 + 2.51 \, x
\tag{N-11}
\end{equation}

$w_0 = 0.90$ = intercept (predicted price when size = 0).

$w_1 = 2.51$ = slope (each extra sqm adds 2.51 lakh to the price).


\subsection*{Step 7: Verify --- Do Predictions Match the Data?}

With only 2 data points and 2 parameters ($w_0, w_1$), the line must pass \emph{exactly} through both points:

\begin{align}
\hat{y}_1 &= 0.90 + 2.51 \times 30 = 0.90 + 75.30 = 76.20 \;\; \checkmark \tag{N-12a} \\
\hat{y}_2 &= 0.90 + 2.51 \times 80 = 0.90 + 200.80 = 201.70 \;\; \checkmark \tag{N-12b}
\end{align}

\gut{
With 2 points and 2 unknowns, there is exactly one line through both points --- zero residual error. This is like solving 2 equations with 2 unknowns in high school algebra. It's only when you have more data points than parameters (overdetermined system) that you get nonzero residuals, and the formula finds the \emph{best compromise} line.
}


\subsection*{Summary}

\begin{center}
\renewcommand{\arraystretch}{1.3}
\begin{tabular}{ccc}
\toprule
\textbf{Step} & \textbf{Operation} & \textbf{Result} \\
\midrule
1 & Build $X$, $\mathbf{y}$ & $2\times2$ matrix, $2\times1$ vector \\
2 & Transpose $X$ & $X^\top$ ($2\times2$) \\
3 & $X^\top X$ & $\begin{pmatrix} 2 & 110 \\ 110 & 7300 \end{pmatrix}$ \\
4 & $X^\top \mathbf{y}$ & $\begin{pmatrix} 277.9 \\ 18422 \end{pmatrix}$ \\
5 & $(X^\top X)^{-1}$ & $\begin{pmatrix} 2.92 & -0.044 \\ -0.044 & 0.0008 \end{pmatrix}$ \\
6 & $(X^\top X)^{-1} X^\top \mathbf{y}$ & $\begin{pmatrix} 0.90 \\ 2.51 \end{pmatrix}$ \\
\bottomrule
\end{tabular}
\end{center}

\end{document}
